\section{Riferimenti}
\label{sec:riferimenti}

\begin{enumerate}[leftmargin=1.4cm, label={[\arabic*]}]
    \item \label{bib:corso} Materiale didattico del corso di \emph{Internet of Things},
    Prof. Silvio Barra, Corso di Laurea Magistrale in Informatica,
    Università degli Studi di Napoli Federico II, A.A. 2025/2026.

    \item \label{bib:dataset} Dataset \emph{PlantVillage -- Tomato Leaf Dataset}, disponibile su
    Kaggle: \url{https://www.kaggle.com/datasets/charuchaudhry/plantvillage-tomato-leaf-dataset}.

    \item \label{bib:cnntomato} Progetto \texttt{CNN-TOMATO-}: confronto sperimentale FCNN vs CNN
    sul riconoscimento di patologie fogliari del pomodoro, progetto svolto
    nell'ambito dell'esame di \emph{Neural Networks} --- fonte del modello
    \texttt{TomatoCNN} riutilizzato in PomodorIA.

    \item \label{bib:pytorch} Documentazione ufficiale di PyTorch: \url{https://pytorch.org/docs/}
    --- in particolare le sezioni relative a \texttt{torch.nn.utils.prune}
    (pruning) e \texttt{torch.quantization} (quantizzazione dinamica).

    \item \label{bib:onnx} Documentazione ufficiale di ONNX e ONNX Runtime:
    \url{https://onnx.ai/} e \url{https://onnxruntime.ai/} --- in
    particolare gli strumenti di esportazione (\texttt{torch.onnx.export})
    e di quantizzazione (\texttt{onnxruntime.quantization}).

    \item \label{bib:streamlit} Documentazione ufficiale di Streamlit: \url{https://docs.streamlit.io/}
    --- framework utilizzato per la realizzazione della dashboard
    interattiva di monitoraggio.

    \item \label{bib:rn} S. Russell, P. Norvig, \emph{Artificial Intelligence: A Modern
    Approach} --- riferimento standard per la formalizzazione degli agenti
    razionali e del modello PEAS.

    \item \label{bib:earlyblight} NC State Extension, \emph{Early Blight of Tomato}:
    \url{https://content.ces.ncsu.edu/early-blight-of-tomato} --- condizioni
    ambientali (temperatura, umidità) favorevoli allo sviluppo dell'Early
    Blight, usate come riferimento qualitativo per calibrare i profili
    ambientali del simulatore.

    \item \label{bib:lateblight} UC IPM (University of California, Statewide Integrated
    Pest Management Program), \emph{Late Blight -- Tomato}:
    \url{https://ipm.ucanr.edu/agriculture/tomato/late-blight/} ---
    condizioni ambientali favorevoli allo sviluppo del Late Blight, usate
    con lo stesso criterio del riferimento precedente.

    \item \label{bib:datafusion} D.L. Hall, J. Llinas, \emph{An Introduction to
    Multisensor Data Fusion}, Proceedings of the IEEE, 85(1), 1997 ---
    riferimento classico sulla combinazione di dati da fonti eterogenee.

    \item \label{bib:fog} F. Bonomi, R. Milito, J. Zhu, S. Addepalli, \emph{Fog
    Computing and its Role in the Internet of Things}, Proceedings of the
    MCC Workshop on Mobile Cloud Computing, ACM, 2012 --- introduce il
    concetto di Fog Computing come estensione del Cloud verso il margine
    della rete.
\end{enumerate}
